\chapter{Încărcarea datelor}
Având în vedere formatele diferite în care informațiile extrase din fișierele audio pot fi salvate pe disc, 
același nivel de flexibilitate este necesar și la încărcarea acestor date în memorie. Nevoia de a
implementa diferite moduri de încărcare pentru tipuri diferite de formate ale datelor a fost un rezultat al capabilităților 
mașinilor pe care rețeaua a fost antrenată. În funcție de aceste capabilități, dar și în funcție de dimensiunea informațiilor 
extrase, datele puteau sau nu să fie încărcate în totalitate în memoria RAM. Această problemă a fost rezolvată prin experimentarea 
cu diferite tipuri de generatoare care încarcă datele de pe disc în timpul antrenării rețelei, scopul final fiind găsirea unei metode care 
minimizează timpul de antrenare.

O schemă generală a încărcării datelor în memorie și pregătirea inputurilor de antrenare, 
validare și test ale rețelei este prezentată în figura \ref{fig:f13}

\begin{figure}[h]
    \caption{Procesul de încărcare al datelor în memorie}
    \includegraphics[width=0.98\textwidth]{data load.png}
    \centering
    \label{fig:f13}
\end{figure} 

După cum a fost menționat, au existat două abordări, în funcție de dimensiunea datelor și memoria RAM disponibilă. 
În continuare sunt descrise particularitățile acestor abordări.

\section{Încărcarea întreg setului de date în memorie}
În acest caz, pipeline-ul primește ca input o instanță a unei clase derivate din clasa de bază Features Loader și, opțional, 
un state random ce va fi utilizat la amestecarea datelor.

Clasa Features Loader se ocupă de încărcarea informațiilor din fișiere și extragerea genurilor muzicale. 
În funcție de clasa concretă specificată, informațiile sunt încărcate din fișiere JSON, deserializate ca Tensori din fișiere binare, 
sau citite ca JSON linie cu linie dintr-un fișier, iar genurile sunt extrase din numele fișierelor, respectiv din obiectul JSON citit.
La instanțiere se specifică calea directorului sau a fișierului în care se află datele și un subset al genurilor din setul de date. 
Opțiunea de a specifica un subset de genuri este utilizată în procesul ulterior de selecție a genurilor muzicale dezvoltat în secțiunea 6.1.6. 
În urma apelării funcției de încărcare se obțin trei liste ce corespund informațiilor extrase din fișierele audio prezente în setul de date, 
genurilor și id-urilor corespunzătoare acestor informații.

Fiecărui gen distinct încărcat îi este atribuit un număr întreg iar fiecare gen din listă este înlocuit cu numărul ce îi corespunde. 
În cazul informațiilor extrase din fișierele audio, dacă mai multe tipuri de informații sunt prezente, acestea vor fi concatenate. 

Setul de date obținut va fi amestecat și împărțit în trei seturi de antrenare, test și validare conform unor procente ce pot fi
specificate. Seturile de date rezultate pot fi 
liste sau Tensorflow Datasets \footnote{\href{https://www.tensorflow.org/api_docs/python/tf/data/Dataset}{tensorflow.data.Dataset}}.

Cu toate că utilizarea Tensorflow Datasets în favoarea listelor rezultă într-o antrenare mai rapidă a rețelei, în funcție de dimensiunea
datelor, nu este posibilă alocarea memoriei pentru crearea obiectelor tensorflow.data.Dataset. În aceste cazuri, au fost folosite listele
ca input pentru rețea sau generatorii.

\section{Generatori}
O soluție a problemei de memorie este utilizarea obiectelor Tensorflow Datasets construite peste generatori. Astfel, setul de date nu 
mai este încărcat în totalitate în memorie. Cand se folosește această abordare pipeline-ul primește un obiect al clasei File Names 
Features Loader, derivată din clasa Features Loader  care încarcă doar numele tuturor fișierelor din setul de date și genurile 
corespunzătoare. Procesul de amestecare, înlocuire a genurilor cu numere întregi și separare în date de antrenare, test și validare 
este același. Diferența constă în faptul că, după separarea datelor, sunt instanțiate trei generatoare ce se vor ocupa cu citirea
informațiilor din fișier. Sunt create trei obiecte Tensorflow Datasets pe baza generatoarelor, utilizând functia $from\_generator$
\footnote{\href{https://www.tensorflow.org/api_docs/python/tf/data/Dataset\#from\_generator}{tensorflow.data.Dataset.from\_generator}}.

Tensorflow datasets se ocupă automat de formarea batch-urilor și preîncărcarea datelor, fiind mai rapidă decât utilizarea generatorilor
creați ca input direct al rețelei. Datorită faptului că utilizarea Tensorflow Datasets îmbunătățește timpul de antrenare al rețelei în comparație
cu listele simple, în cazurile în care dimensiunea memoriei RAM permitea acest lucru, s-au utilizat generatori care încărcau toate datele în
memorie și le returnau pe rând utilizând Tensorflow Datasets. În acest mod, a fost rezolvată problema spațiului, mentionată anterior, 
atunci când se alocă
obiecte Tensorflow Datasets direct pentru listele încărcate în memorie.

În continuare se pot observa rezultatele experimentelor atunci când s-au utilizat modurile descrise de încărcare a datelor.
\section{Experimente}